% !TeX root = ../se200_referential_regimes.tex

% =============================================================================
\section{Profile Distinctions and Carrier-Indexed Combination}
\label{sec:distinctions}
% =============================================================================

The preceding sections specify candidate carrier kinds,
identity bases, and transformation classifications.
This section compares the resulting profiles by their classification
functions and by the relations they generate, and shows how
carrier-specific relations can be combined in one tagged relation.



% -----------------------------------------------------------------------------
\subsection{Witnessed Differences Within the Selected Profiles}
\label{subsec:witnessed_splits}
% -----------------------------------------------------------------------------

The proposed core contains pairs of identity bases associated
with the same carrier kind.
These pairs motivate tests of whether ordinary transformations make
different identity judgments under the respective bases.

For plain referents, the proposed locus-fixed and object-fixed bases
are compared using branch/fork operations
(Proposition~\ref{se200.prop.PlainReferentUnderdetermination}).
A particular fork may preserve a locus while distinguishing the objects
associated with that locus.
The question for a realization is which bearer each basis tracks and
whether the proposed preservation and breakage judgments are warranted.

For scopes, the proposed extension-fixed and structure-fixed bases
are compared using aggregation and decomposition
(Proposition~\ref{se200.prop.ScopeUnderdetermination}).
An operation may preserve coverage while changing component structure.
Whether that difference is an identity difference depends on the
independently justified criterion for the scope in question.

For rules, the proposed content-fixed and structure-fixed bases
are compared using refinement
(Proposition~\ref{se200.prop.RuleUnderdetermination}).
A refinement may preserve specified content while changing structure.
Again, the relevant sameness judgment depends on what the rule is
and on which criterion of its identity is justified.

Two justified identity relations that disagree on an eligible pair
of presentations are distinct on that comparison domain.
A difference between declared transformation classifications does not
by itself exhibit a pair on which the generated relations disagree.
Such a witness requires a shared eligible comparison domain and
appropriate transformation and identity interpretations.
These pairs supply candidate distinctions for further testing.

% -----------------------------------------------------------------------------
\subsection{Profiles and Generated Relations}
\label{subsec:profiles_vs_relations}
% -----------------------------------------------------------------------------

The profiles catalogued in Section~\ref{sec:profiles}
arise from the selected carrier--basis assignments.
Proposition~\ref{se200.prop.DistinctCoreProfiles} concerns distinctions
among their classification functions.

The number of distinct family-level generated relations
depends on the generating rule.
The \PRS{}-only and \PRS{}-or-\NEU{} constructions
induce different patterns of equal family sets;
Appendix~\ref{app:sep} enumerates them and gives the resulting
upper bounds.
Further equalities may arise when the dynamics do not distinguish
the selected family sets.
These bounds do not apply to the invariant-based step relations,
whose equality and separation are characterized by
Theorem~\ref{se200.thm.StepIdentityCollapse} and
Proposition~\ref{se200.prop.StepSeparation}.

% -----------------------------------------------------------------------------
\subsection{Carrier-Indexed Combination}
\label{subsec:amalgamation}
% -----------------------------------------------------------------------------

An elementary construction shows that relations on distinct carrier
kinds can be represented together while preserving each component
equivalence relation.
Let $K$ index the selected carrier kinds, and let $D_k$ be a domain
of presentations for carrier kind $k$.
Suppose $E_k$ is an equivalence relation on $D_k$ for each $k\in K$.
Use the tagged disjoint union
\[
D = \bigsqcup_{k\in K} D_k
\]
and define a relation $E$ on $D$ by
\[
(k,x)\mathrel{E}(\ell,y)
\quad\Longleftrightarrow\quad
\bigl(k=\ell \;\wedge\; x\mathrel{E_k} y\bigr).
\]

\begin{proposition}[Carrier-indexed combination]
\label{se200.prop.CarrierIndexedCombination}
Then $E$ is an equivalence relation on $D$.
Restricting $E$ to the $k$-th summand recovers $E_k$
for each $k\in K$.
\end{proposition}

\begin{proof}
Reflexivity and symmetry follow from the corresponding properties
of each $E_k$.
If $(k,x)\mathrel{E}(\ell,y)$ and
$(\ell,y)\mathrel{E}(m,z)$, then $k=\ell=m$ and
$x\mathrel{E_k}y\mathrel{E_k}z$.
Transitivity of $E_k$ therefore gives $(k,x)\mathrel{E}(m,z)$.
For presentations in the $k$-th summand, the definition of $E$
is exactly the definition of $E_k$.
\end{proof}

This construction preserves the component equivalence relations
without identifying presentations belonging to different carrier kinds.
Distinct carrier kinds therefore do not by themselves require
separate formal structures: a declared carrier tag carries the
distinction within one relation.
Whether a combined relation is an adequate identity account, preserves
all required transformation semantics, and satisfies the Neutral
Substrate constraint depends on the semantics of the tags and on the
bases proposed for each carrier kind.
A publicly declared carrier or basis parameter is not a hidden regime
merely because it is a parameter; its adequacy depends on whether its
semantics are fixed and its use is justified.

% -----------------------------------------------------------------------------
\subsection{Separation Under an Extended Transformation Basis}
\label{subsec:monotonicity}
% -----------------------------------------------------------------------------

A narrower monotonicity result applies to the declared
classification profiles.
Let $\Fset\subseteq\Fset'$ be transformation bases.
Suppose the classifications of two profiles on $\Fset'$ agree
with their existing classifications on every operation in $\Fset$.

\begin{proposition}[Persistence of a classification distinction]
\label{se200.prop.ClassificationDistinctionPersists}
If two profiles have different classifications for some
$f\in\Fset$, their classification functions remain distinct
under extensions to $\Fset'$ that preserve the original classifications.
\end{proposition}

\begin{proof}
The extended classification functions still disagree at $f$,
so they cannot be equal.
\end{proof}

This result records that a classification difference, once established,
is stable under extensions of the transformation basis that retain the
original classifications.

The broader limitations are discussed in Section~\ref{sec:limits}.
